% ============================================================================
%  SPEAR3 RF System documentation -- shared preamble
%  Engine: LuaLaTeX.  Build with build.ps1 (latexmk -lualatex).
% ============================================================================

\usepackage[letterpaper,top=1in,bottom=1in,left=1.05in,right=1.05in]{geometry}

% --- fonts -----------------------------------------------------------------
\usepackage{fontspec}
\setmainfont{TeX Gyre Termes}
\setsansfont{TeX Gyre Heros}
\setmonofont{Consolas}[Scale=0.86]

% --- Unicode safety --------------------------------------------------------
% LuaLaTeX DROPS unmapped glyphs SILENTLY -- an unmapped U+03BC turns "8 uF"
% into "8 F".  Every non-ASCII character used in these documents is mapped
% here.  The build script fails if the log reports "Missing character".
\usepackage{newunicodechar}
\newunicodechar{μ}{\ensuremath{\mu}}      \newunicodechar{µ}{\ensuremath{\mu}}
\newunicodechar{Ω}{\ensuremath{\Omega}}   \newunicodechar{≈}{\ensuremath{\approx}}
\newunicodechar{≥}{\ensuremath{\geq}}     \newunicodechar{≤}{\ensuremath{\leq}}
\newunicodechar{≠}{\ensuremath{\neq}}     \newunicodechar{×}{\ensuremath{\times}}
\newunicodechar{÷}{\ensuremath{\div}}     \newunicodechar{±}{\ensuremath{\pm}}
\newunicodechar{−}{\ensuremath{-}}        \newunicodechar{√}{\ensuremath{\sqrt{}}}
\newunicodechar{α}{\ensuremath{\alpha}}   \newunicodechar{β}{\ensuremath{\beta}}
\newunicodechar{φ}{\ensuremath{\varphi}}  \newunicodechar{ω}{\ensuremath{\omega}}
\newunicodechar{Δ}{\ensuremath{\Delta}}   \newunicodechar{τ}{\ensuremath{\tau}}
\newunicodechar{θ}{\ensuremath{\theta}}   \newunicodechar{λ}{\ensuremath{\lambda}}
\newunicodechar{π}{\ensuremath{\pi}}      \newunicodechar{σ}{\ensuremath{\sigma}}
\newunicodechar{₀}{\ensuremath{_0}}       \newunicodechar{₁}{\ensuremath{_1}}
\newunicodechar{₂}{\ensuremath{_2}}       \newunicodechar{²}{\ensuremath{^2}}
\newunicodechar{³}{\ensuremath{^3}}       \newunicodechar{⁴}{\ensuremath{^4}}
\newunicodechar{⁻}{\ensuremath{^-}}       \newunicodechar{⁺}{\ensuremath{^+}}
\newunicodechar{→}{\ensuremath{\rightarrow}}
\newunicodechar{←}{\ensuremath{\leftarrow}}
\newunicodechar{↔}{\ensuremath{\leftrightarrow}}
\newunicodechar{⇒}{\ensuremath{\Rightarrow}}
\newunicodechar{⊕}{\ensuremath{\oplus}}   \newunicodechar{∥}{\ensuremath{\parallel}}
\newunicodechar{°}{\ensuremath{^\circ}}   \newunicodechar{·}{\ensuremath{\cdot}}
\newunicodechar{≪}{\ensuremath{\ll}}      \newunicodechar{≫}{\ensuremath{\gg}}
\newunicodechar{✓}{\checkmark}            \newunicodechar{⚠}{\textbf{!}}
\newunicodechar{–}{--}                    \newunicodechar{—}{---}

% --- maths and units -------------------------------------------------------
\usepackage{amsmath,amssymb}
\usepackage{siunitx}
\DeclareSIUnit\voltampere{VA}
\DeclareSIUnit\VAR{var}
\sisetup{per-mode=symbol,range-phrase={ to },range-units=single,
         group-minimum-digits=4,detect-weight=true}

% --- tables ----------------------------------------------------------------
\usepackage{longtable,booktabs,array,tabularx,multirow}
\usepackage{ragged2e}
\setlength{\LTleft}{0pt}\setlength{\LTright}{0pt}
\renewcommand{\arraystretch}{1.18}
% left-aligned, hyphenating paragraph column used by every wide table
\newcolumntype{P}[1]{>{\RaggedRight\arraybackslash}p{#1}}

% --- graphics and figures --------------------------------------------------
\usepackage{graphicx}
\graphicspath{{figures/}{figures/photos/}{figures/generated/}}
\usepackage{float}
\usepackage[font=small,labelfont=bf,width=0.92\linewidth]{caption}
\usepackage{comment}

% A figure must not drift out of the section it belongs to.  With 47 floats
% LaTeX will otherwise push them pages away from the text that discusses them.
\usepackage[section]{placeins}

% Default float parameters are tuned for papers with two or three figures.
% This document is figure-dense, so allow more of a page to be float and more
% floats per page; without this, figures queue up and drift several pages.
\renewcommand{\topfraction}{0.88}
\renewcommand{\bottomfraction}{0.70}
\renewcommand{\textfraction}{0.10}
\renewcommand{\floatpagefraction}{0.72}
\setcounter{topnumber}{3}
\setcounter{bottomnumber}{2}
\setcounter{totalnumber}{5}

% Image size policy.  Photographs are capped in BOTH width and height so that
% no figure can ever exceed roughly half a page, whatever its aspect ratio.
\newlength{\photomaxh}\setlength{\photomaxh}{0.34\textheight} % ~1/3 page
\newlength{\widemaxh} \setlength{\widemaxh}{0.46\textheight}  % ~1/2 page

% \photofig{file}{caption}{label}      -- default, about 1/3 page
\newcommand{\photofig}[3]{%
  \begin{figure}[htbp]\centering
    \includegraphics[width=0.62\linewidth,height=\photomaxh,keepaspectratio]{#1}%
    \caption{#2}\label{#3}
  \end{figure}}
% \widefig{file}{caption}{label}       -- for detailed schematics, about 1/2 page
\newcommand{\widefig}[3]{%
  \begin{figure}[htbp]\centering
    \includegraphics[width=0.94\linewidth,height=\widemaxh,keepaspectratio]{#1}%
    \caption{#2}\label{#3}
  \end{figure}}
% \smallfig{file}{caption}{label}      -- for low-resolution sources
\newcommand{\smallfig}[3]{%
  \begin{figure}[htbp]\centering
    \includegraphics[width=0.44\linewidth,height=0.26\textheight,keepaspectratio]{#1}%
    \caption{#2}\label{#3}
  \end{figure}}
% \tallfig{file}{caption}{label}       -- portrait-orientation photographs.
% A tall photo capped at \photomaxh comes out only ~45 mm wide, so these are
% capped by height instead and allowed a narrower column.
\newcommand{\tallfig}[3]{%
  \begin{figure}[tbp]\centering
    \includegraphics[width=0.46\linewidth,height=0.50\textheight,keepaspectratio]{#1}%
    \caption{#2}\label{#3}
  \end{figure}}
% two portrait photographs side by side under one caption:
% \tallpairfig{fileA}{subA}{labA}{fileB}{subB}{labB}{caption}{label}
\newcommand{\tallpairfig}[8]{%
  \begin{figure}[tbp]\centering
    \begin{minipage}[b]{0.46\linewidth}\centering
      \includegraphics[width=\linewidth,height=0.40\textheight,keepaspectratio]{#1}%
      \subcaption{#2}\label{#3}
    \end{minipage}\hfill
    \begin{minipage}[b]{0.46\linewidth}\centering
      \includegraphics[width=\linewidth,height=0.40\textheight,keepaspectratio]{#4}%
      \subcaption{#5}\label{#6}
    \end{minipage}
    \caption{#7}\label{#8}
  \end{figure}}
% two photographs side by side, each about 1/4 page
\newcommand{\pairfig}[6]{%
  \begin{figure}[htbp]\centering
    \begin{minipage}[t]{0.48\linewidth}\centering
      \includegraphics[width=\linewidth,height=0.26\textheight,keepaspectratio]{#1}%
      \subcaption{#2}\label{#3}
    \end{minipage}\hfill
    \begin{minipage}[t]{0.48\linewidth}\centering
      \includegraphics[width=\linewidth,height=0.26\textheight,keepaspectratio]{#4}%
      \subcaption{#5}\label{#6}
    \end{minipage}
  \end{figure}}
\usepackage{subcaption}

% --- diagrams --------------------------------------------------------------
\usepackage{tikz}
\usetikzlibrary{arrows.meta,positioning,fit,backgrounds,calc,shapes.geometric,
                shapes.misc,chains,decorations.pathreplacing,patterns,matrix}
\usepackage{pgfplots}\pgfplotsset{compat=1.18}
\usepackage{rotating}                  % sidewaysfigure, if a portrait page is wanted
\usepackage{pdflscape}                 % landscape: rotates the PAGE in the PDF
\usepackage[export]{adjustbox}         % max width= / max height= on a box

% Shrink a diagram to the text width only if it would otherwise overflow, so
% figures drawn at their natural scale are never enlarged or needlessly shrunk.
\newsavebox{\fitbox}
\newenvironment{fitpicture}
  {\begin{lrbox}{\fitbox}}
  {\end{lrbox}%
   \ifdim\wd\fitbox>\linewidth
     \resizebox{\linewidth}{!}{\usebox{\fitbox}}%
   \else
     \usebox{\fitbox}%
   \fi}

% House diagram style.  Keep every block diagram in these styles so the
% document reads as one drawing set rather than a scrapbook.
\definecolor{cCtrl}{HTML}{1F4E79}   % control / EPICS
\definecolor{cPlc} {HTML}{7F6000}   % PLC / interlock
\definecolor{cPwr} {HTML}{843C0C}   % high power
\definecolor{cRf}  {HTML}{2E6B34}   % RF signal path
\definecolor{cNeut}{HTML}{3B3B3B}
\tikzset{
  blk/.style   ={draw=cNeut,fill=white,rounded corners=1.6pt,align=center,
                 inner sep=4pt,font=\scriptsize,minimum height=8mm},
  ctrl/.style  ={blk,draw=cCtrl,fill=cCtrl!7,text=cCtrl!85!black},
  plc/.style   ={blk,draw=cPlc, fill=cPlc!8, text=cPlc!80!black},
  pwr/.style   ={blk,draw=cPwr, fill=cPwr!7, text=cPwr!85!black},
  rf/.style    ={blk,draw=cRf,  fill=cRf!7,  text=cRf!80!black},
  grp/.style   ={draw=black!35,dashed,rounded corners=3pt,inner sep=7pt},
  lbl/.style   ={font=\scriptsize\itshape,text=black!62},
  sig/.style   ={-{Stealth[length=2.2mm]},draw=cNeut,thick},
  fib/.style   ={sig,densely dashed},
  bus/.style   ={draw=cNeut,line width=1.1pt},
  tag/.style   ={font=\tiny,inner sep=1.4pt,fill=white,text=black!70},
}
% consistent caption for a TikZ figure
\newcommand{\diagcap}[2]{\caption{#1}\label{#2}}

% --- lists, misc -----------------------------------------------------------
\usepackage{enumitem}
\setlist{nosep,leftmargin=1.35em}
\providecommand{\tightlist}{%   emitted by pandoc for compact lists
  \setlength{\itemsep}{0pt}\setlength{\parskip}{0pt}}
\usepackage[table]{xcolor}
\usepackage{soul}
\usepackage{microtype}
\usepackage{fancyhdr}
\usepackage{lastpage}
\usepackage{needspace}

% --- callout box for warnings / key findings -------------------------------
\usepackage[most]{tcolorbox}
\newtcolorbox{keynote}[1][]{colback=black!3,colframe=black!45,boxrule=0.4pt,
  arc=1.5pt,left=5pt,right=5pt,top=4pt,bottom=4pt,fontupper=\small,#1}
\newtcolorbox{warnbox}[1][]{colback=cPwr!5,colframe=cPwr!70,boxrule=0.7pt,
  arc=1.5pt,left=5pt,right=5pt,top=4pt,bottom=4pt,fontupper=\small,#1}

% --- cross references and links --------------------------------------------
\usepackage[hyphens]{url}
\usepackage[hidelinks,pdfusetitle,bookmarksnumbered]{hyperref}
\usepackage[capitalise,nameinlink]{cleveref}
\crefname{section}{\S}{\S\S}
\Crefname{section}{Section}{Sections}

% Repository paths are long and unhyphenatable; allow breaks after / and .
% Long file paths and CamelCase document names have no natural break points, so
% url's default break set is not enough: allow a break before any capital and
% after the usual separators.
\newcommand{\fpath}[1]{{\ttfamily\small\Urlmuskip=0mu plus 1mu\relax
  \def\UrlBreaks{\do\/\do\.\do\-\do\_\do\:\do\A\do\B\do\C\do\D\do\E\do\F\do\G
    \do\H\do\I\do\J\do\K\do\L\do\M\do\N\do\O\do\P\do\Q\do\R\do\S\do\T\do\U
    \do\V\do\W\do\X\do\Y\do\Z}%
  \urlstyle{tt}\nolinkurl{#1}}}
\setlength{\emergencystretch}{2em}

% --- headers ---------------------------------------------------------------
\pagestyle{fancy}\fancyhf{}
\renewcommand{\headrulewidth}{0.3pt}
\fancyhead[L]{\footnotesize\nouppercase{\rightmark}}
\fancyhead[R]{\footnotesize\thepage\ / \pageref{LastPage}}
\renewcommand{\sectionmark}[1]{\markright{\thesection\quad #1}}
\fancypagestyle{plain}{\fancyhf{}\renewcommand{\headrulewidth}{0pt}%
  \fancyfoot[C]{\footnotesize\thepage}}

\setcounter{tocdepth}{2}
\setcounter{secnumdepth}{3}
